\section{Introduction}\label{sec:intro}

On June 11, 2026, eight European researchers, investors, and policy writers published \emph{Europe 2031}, a scenario of how the continent loses its place in the age of artificial intelligence \citep{juijn2026europe}. Its climax is a phone call. In March 2031 the US vice president proposes to the leaders of the Netherlands, Germany, and France that ASML, the Dutch firm that alone builds the extreme ultraviolet (EUV) lithography machines through which every leading-edge AI chip passes, be folded into a joint Dutch-American holding company whose board gives Washington a controlling vote over output, customers, and technology transfer. When they decline, the offer becomes an ultimatum. The authors call ASML ``the only card Europe still holds'' and ``the one bottleneck the entire AI race runs through''. Yet in their story Europe never plays it. In 2028 the Dutch halt ASML's remaining exports and servicing of older immersion lithography tools in China after Washington threatens to apply its foreign direct product rule to the firm, and Europe ``negotiates nothing in return''. In 2029 a package of EU countermeasures under the Anti-Coercion Instrument, including restrictions on ASML's exports and servicing to American fabs in Arizona, falls short of a qualified majority. In 2031 Europe is left choosing among submission to Washington, alignment with Beijing, and isolation.

The puzzle is not confined to fiction. Since 2019 the Netherlands has aligned its export controls with American demands in a series of steps: it let ASML's licence to ship an EUV machine to China lapse without renewal in 2019, imposed national licensing on advanced immersion lithography in 2023, partly revoked a licence in January 2024, and extended its controls to older immersion tools and their spare parts and software in September 2024 \citep{nltimes2020pressured, govnl2023export, dutchnews2024licence, asml2024license}. The main thing The Hague received in return was an exemption, in December 2024, from American rules that would otherwise have applied to the same tools, and Washington can withdraw it by amending a rule \citep{bis2024fdpr}. When members of Congress moved in 2026 to ban sales of immersion lithography tools anywhere in China and the servicing of tools at China's leading fabs, the Dutch trade minister protested that the bill ``could potentially put our national security in American hands'' \citep{nltimes2026irritated}. If ASML is the chokepoint of the most consequential technology race of the century, why does its owner seem to extract so little for it?

The literatures that should answer this question point the other way. Theories of asymmetric interdependence hold that the less dependent party to an exchange gains influence over the more dependent one \citep{hirschman1945national, keohane1977power}. Work on weaponized interdependence argues that states with jurisdiction over central nodes in global networks can turn them into instruments of coercion \citep{farrell2019weaponized}. Both imply that the United States, which cannot extend its lead in AI without ASML's machines, should be paying Europe for continued access. Formal models of sanctions and issue linkage analyze a sender and a target \citep{eaton1992sanctions, lacy2004theory}. Work on alliance coercion studies a patron pressing a dependent client \citep{gerzhoy2015alliance}, and work on coercive assurance shows that what makes a threat credible can make the accompanying promise less so \citep{pauly2024damned}. Recent work in geoeconomics locates a hegemon's power in its ability to coordinate threats across relationships \citep{clayton2026framework}. None of these considers a dependent ally that holds a chokepoint its patron needs, in a race in which a rival great power would pay for access.

We develop a game-theoretic model of that situation. Europe controls a chokepoint and decides whether to keep it aligned with its security patron, the United States, or to let it serve the patron's rival, China. Both great powers make offers, and each controls instruments linked to Europe's choice: security protection and access to frontier AI on the American side, critical inputs on the Chinese side. A threat to use any of these instruments counts only if carrying it out would be sequentially rational.

We argue that a chokepoint is worth what its holder could credibly gain by walking away, not what its patron would lose without it. We call the first quantity Europe's \emph{exit premium}: the most China can credibly offer, net of the costs the United States can credibly impose and the value Europe itself places on alignment. When the exit premium is negative, Europe receives nothing however much Washington depends on ASML. The patron's dependence returns only through four channels: sanctions whose use would harm the patron itself, the holder's ability to commit to a demand, the risk that a dispute slips out of control, and the holder's ability to inflict pain while bargaining continues.

The model yields five sets of results. First, leverage equals the positive part of the exit premium, and a positive premium is necessary for any leverage under every bargaining protocol we consider in which only Europe's own choice can produce exit. Second, Europe's dependence on American protection gives Washington leverage only to the extent that the guarantee is \emph{conditional}, reliable for a loyal Europe and void for a realigned one. Compute scarcity first creates an American AI lever and then destroys it, once Washington rations loyal allies too. China's capacity to coerce Europe raises what Washington must pay for European alignment but never raises Europe's own payoff, while China's promises help Europe only when it already has leverage. Third, because Europe cannot commit not to exploit a rising exit premium, Washington has a preventive motive to demand control of the chokepoint while its own threats still bite. It needs a public commitment to do so, and a Chinese threat to punish the transfer of control protects European ownership. Fourth, Europe's chokepoint is nested in American ones. ASML depends on US-origin subsystems, US staff, and US customers, so Europe's exit threat is only as credible as the Dutch state's willingness to compel its own champion. Washington can also target the member states whose consent exit requires, although a rival bidder in the Council would force it to pay more than the cost of the cheapest blocking minority. Fifth, the one place where American dependence on ASML matters directly is a war of attrition over servicing. There, what decides the outcome is time to pain: the length of the US stockpile relative to Europe's endurance.

We use the model to reread both the scenario and the record of 2019--2026. The 2028 capitulation in \emph{Europe 2031} is an equilibrium of the game the authors describe. The 2029 vote is not, because the states the scenario names cannot form a blocking minority under 2026 Council population weights. The 2031 ultimatum is a preventive demand for control rights, but its threat of permanent exclusion from American AI is not credible. In the real record, Dutch alignment without compensation, the December 2024 US exemption of the Netherlands and other listed allies from new foreign direct product rules, the Nexperia dispute of 2025, and the 2026 fight over servicing restrictions are what the model predicts for a holder whose exit premium is negative.

The paper contributes to four literatures. First, it qualifies the theory of weaponized interdependence by showing that control of a hub yields leverage over the hub's main user only through a credible alternative, and by formalizing the jurisdictional condition that a nested chokepoint fails. Second, it separates two things the alliance literature often treats together, an ally's dependence on its patron and the conditionality of the patron's commitment, and shows that only the second is a source of the patron's leverage. Third, it extends the analysis of commitment problems under shifting power from preventive war to preventive demands for control rights over an asset that shapes future bargaining power \citep{fearon1995rationalist, fearon1996bargaining, powell2006war}. Fourth, it gives the European debate about technological sovereignty a set of conditions under which ASML is, and is not, a card worth holding.

\Cref{sec:background} describes ASML's position and the three dependencies that shape Europe's use of it. \Cref{sec:literature} sets out the arguments we build on and against. \Cref{sec:model} presents the model and its main results, \cref{sec:control} the analysis of control rights, and \cref{sec:extensions} three extensions. \Cref{sec:revisited} applies the results to the scenario and to the record since 2019. \Cref{sec:implications} draws out implications and limitations.
